\begin{frame}{Why Chain-of-Thought is not enough?}
\small



\begin{block}{CoT improves performance over direct answering. But \dots}
    \begin{itemize}
    \item biased prompts still influence the LLMs output. Up to 36\% precision loss.
    \item even with unbiased prompts, the model reaches it's limitations.
    \item problem: Especially in BBH-Benchmarks, the model still must do two things at once:
    \begin{enumerate}
        \item decompose the problem into useful steps,
        \item execute these steps correctly.
    \end{enumerate}
\end{itemize}
\end{block}

\small

\begin{columns}[T]
\setbeamercolor{block title}{fg=black,bg=blue!20}
% =====================================================
\begin{column}{0.48\textwidth}

\begin{block}{Input}

\textbf{Q:} The bakers at the Beverly Hills Bakery baked (\small\textbf{+})200 loaves of bread on Monday morning. They sold (\small\textbf{-})93 loaves in the morning and (\small\textbf{-})39 loaves in the afternoon. A grocery store returned (\small\textbf{+})6 unsold loaves. How many loaves of bread did they have left?

\end{block}

\end{column}

% =====================================================
\begin{column}{0.48\textwidth}
\setbeamercolor{block title}{fg=black,bg=green!20}
\begin{block}{\CoT{} Output}

\textbf{A:}
The bakers started with 200 loaves. They sold 93 in the morning and 39 in the afternoon. So they sold \(93 + 39 = 132\) loaves. The grocery store returned 6 loaves. So they had \(200 - 132 - 6 = 62\) loaves left.
\[
    \text{{\color{red}$\times$} \underline{Wrong}! }\ 200 - (93 + 39) + 6 = 74 
\] 

\end{block}

\end{column}

% =====================================================
\end{columns}



%\begin{alertblock}{Where the limitation appears}
%\begin{itemize}
%    \item At a certain amount of complexity, e.g. number of objects involved into the problem, the precision drops sharply
%    \item When calculating with large numbers the LLM often gets wrong results
%    \item The reasoning trace may look plausible, but the final answer can still be wrong because the model makes arithmetic, logical, or bookkeeping errors.
%\end{itemize}
%\end{alertblock}

\end{frame}